\documentclass[fleqn, a4paper, 12pt]{article}
\usepackage{amsmath, amssymb, amsthm}
\usepackage{marginnote}
\usepackage{gensymb}
\usepackage{commath}
\usepackage{xcolor}
\usepackage{cancel}
\usepackage{siunitx}
\usepackage{tikz, pgfplots}
	\usetikzlibrary{calc, hobby, patterns, intersections}
\usepackage{graphicx}
\usepackage{hyperref}
\usepackage{datetime}
\usepackage{changes}
\usepackage{ulem}
\usepackage{xfrac}
\usepackage{asymptote}
\usepackage{enumerate}
\usepackage{todonotes}
\usepackage{float}
\setcounter{secnumdepth}{4}
\newcommand\numberthis{\addtocounter{equation}{1}\tag{\theequation}}

\newcommand{\AxisRotator}[1][rotate=0]{%
	\tikz [x=0.25cm,y=0.60cm,line width=.2ex,-stealth,#1] \draw (0,0) arc (-150:150:1 and 1);%
}

\theoremstyle{definition}
\newtheorem{example}{Example}
\newtheorem{definition}{Definition}

\theoremstyle{theorem}
\newtheorem{theorem}{Theorem}

\newenvironment{solution}
{\begin{proof}[Solution]\let\qed\relax}
	{\end{proof}}

\newcommand{\curl}{\mathrm{curl\,}}

%\renewcommand{\int_{min}^{max}}{\int\displaylimits_{min}^{max}}

%opening
\title{Lecture 25}
\author{Aakash Jog}
\date{\formatdate{22}{1}{2015}}

\begin{document}

\maketitle
%\setlength{\mathindent}{0pt}

\tableofcontents

\newpage
\section{}

\begin{example}
	A disk with mass $m_1$ and radius $R$ is rotating with $\omega$. A rod is fixed at its centre and two masses of mass $m_2$ are fixed to the rod with springs of coefficient $k$. The natural length of the springs is $r_0$.\\
	At some $t > 0$, and a small enough $\omega$, the masses are in radial equilibrium. Find the $r$ in such a case.\\
	At some $\Omega$, $r = R$. Find $\Omega$.\\
	In the external frame of reference, find the total mechanical energy.
\end{example}

\begin{solution}
	In the wedge frame of reference,
	\begin{align*}
		k (r - r_0) &= m_2 \omega^2 r\\
		\therefore k r - k r_0 &= m_2 \omega^2 r\\
		\therefore r (k - m_2 \omega^2) &= k r_0\\
		\therefore r &= \dfrac{k r_0}{k - m_2 \omega^2}
	\end{align*}
	~\\
	\begin{align*}
		\therefore R &= \dfrac{k r_0}{k - m_2 \Omega^2}\\
		\therefore k - m_2 \Omega^2 &= \dfrac{k r_0}{R}\\
		\therefore m_2 \Omega^2 &= k - \dfrac{k r_0}{R}\\
		\therefore \Omega^2 &= \dfrac{k}{m_2} \left( 1 - \dfrac{r_0}{R} \right)
	\end{align*}
	~\\
	\begin{align*}
		E &= \dfrac{1}{2} I \omega^2 + 2 \cdot \dfrac{1}{2} k (r - r_0)^2\\
		&= \dfrac{1}{2} \left( \dfrac{1}{2} m_1 R^2 + 2 m_2 r^2 \right) \omega^2 + k (r - r_0)^2\\
		&= \dfrac{1}{2} \left( \dfrac{1}{2} m_1 R^2 + 2 m_2 r^2 \right) \omega^2 + k (r - r_0)^2
	\end{align*}
\end{solution}

\begin{example}
	A particle of mass $m_1$ hits a rod of length $L$ and mass $m_2$, at height $h$ from the rod's centre of mass.\\
	\begin{figure}[H]
		\begin{tikzpicture}
			\def\L{4};
			\def\h{1};
			\def\F{0.5};
			
			\draw [ultra thick] (0,{-\L/2}) -- (0,{\L/2});
			
			\fill (-\h,-\h) circle [radius = 2pt];
			
			\draw [-stealth] (-\h,-\h) -- ++(0:\F);
			
			\begin{scope}[dashed]
				\draw ({-\L},0) -- ({\L,0});
				\draw ({-\L},{-\h}) -- ({\L},{-\h});
			\end{scope}
			
			\begin{scope}
				\draw [|<->|] [xshift = 10] (0,{-\L/2}) -- (0,{\L/2}) node [midway, fill = white] {$L$};
				\draw [|<->|] [xshift = 20] (0,0) -- (0,-\h) node [midway, fill = white] {$h$};
			\end{scope}
		\end{tikzpicture}
	\end{figure}
	Assuming that the collision is elastic, find $h$ so that $m_1$ stops after the collision. Find the range of $\dfrac{m_2}{m_1}$ such that $m_1$ stops after the collision.
\end{example}

\begin{solution}
	As the collision is elastic,
	\begin{align*}
		m_1 v_0 &= m_2 v\\
		m_1 v_0 h &= \dfrac{m_2 L^2}{12} \omega\\
		\dfrac{1}{2} m_1 {v_0}^2 &= \dfrac{1}{2} \dfrac{m_2 L^2}{12} \omega^2 + \dfrac{1}{2} m_2 v^2
	\end{align*}
	Therefore,
	\begin{align*}
		m_2 v h &= \dfrac{m_2 L^2}{12} \omega\\
		\therefore v h &= \dfrac{L^2}{12} \omega\\
		\therefore \omega &= \dfrac{12 v h}{L^2}\\
		\therefore \dfrac{1}{2} m_1 {v_0}^2 &= \dfrac{m_2 L^2}{24} \cdot \dfrac{144 v^2 h^2}{L^4} + \dfrac{1}{2} m_2 v^2\\
		\therefore h &= \sqrt{\dfrac{1}{2} L^2 \left( \dfrac{m_2}{m_1} - 1 \right)}
	\end{align*}
	This condition holds if and only if $h \leq \dfrac{L}{2}$.
	\begin{align*}
		0 \leq \dfrac{L}{\sqrt{12}} \sqrt{\dfrac{m_2}{m_1} - 1} \leq \dfrac{L}{2}\\
		1 \leq \dfrac{m_2}{m_1} \leq 4
	\end{align*}
\end{solution}

\begin{example}
	A body made of a hemisphere of radius $a$ and a cylinder of radius $a$ and height $h$ has mass $m$. $\theta$ is the angle between the axis of the vertical. Let the gravitational potential energy be zero at the ground level. Given $z_{\textnormal{COM}} = \left( a + \dfrac{\dfrac{h^2}{2} - \dfrac{a^2}{4}}{h + \dfrac{2a}{3}} \cos \theta \right)$ with respect to the centre of the hemisphere, on the cylinder side of the body, find the gravitational potential energy as a function of $\theta$. Find the condition for the equilibrium at $\theta = 0$ to be stable.
	\begin{figure}[H]
		\begin{tikzpicture}
			\def\R{4};
			\def\h{4};
			\def\angle{30};
			
			\draw ({180 - \angle}:\R) arc ({180 -\angle}:{360 - \angle}:\R);
			
			\draw ({180 -\angle}:\R) --++({90 - \angle}:\h) -- ++({-\angle}:{2*\R}) -- ++({-90 - \angle}:\R) -- cycle;
			
			\draw (-\R,-\R) -- (\R,-\R);
		\end{tikzpicture}
	\end{figure}
\end{example}

\begin{solution}
	\begin{align*}
		y_{\textnormal{COM}} &= a + z_{\textnormal{COM}} \cos \theta\\
		\therefore U &= m g \left( a + \dfrac{\dfrac{h^2}{2} - \dfrac{a^2}{4}}{h + \dfrac{2a}{3}} \cos \theta \right)\\
		\therefore U' &= m g \left( a + \dfrac{\dfrac{h^2}{2} - \dfrac{a^2}{4}}{h + \dfrac{2a}{3}} \cos \theta \right) (- \sin \theta)\\
		\therefore U'' &= - m g \left( a + \dfrac{\dfrac{h^2}{2} - \dfrac{a^2}{4}}{h + \dfrac{2a}{3}} \cos \theta \right)
	\end{align*}
	For the equilibrium at $\theta = 0$ to be stable,
	\begin{align*}
		U'' > 0\\
		\therefore \left( a + \dfrac{\dfrac{h^2}{2} - \dfrac{a^2}{4}}{h + \dfrac{2a}{3}} \cos \theta \right) < 0
	\end{align*}
\end{solution}

\begin{example}
	A solid hemisphere of mass $m$ and radius $a$ is pivoted at its centre, and is used as a pendulum. Find the angular frequency of its oscillations.
\end{example}

\begin{solution}
	\begin{align*}
		\omega &= \sqrt{\dfrac{d m g}{I_0}}\\
		&= \sqrt{\dfrac{\left( \sfrac{3}{8} \right) a m g}{\left( \sfrac{2}{5} \right) m a^2}}\\
		&= \sqrt{\dfrac{15 g}{16 a}}
	\end{align*}
\end{solution}

\begin{example}
	A system is arranged on the ground as shown.
	\begin{figure}[H]
		\begin{tikzpicture}
			\def\L{2};
			
			\draw [gray] ({-2*\L},0) -- ({2*\L},0);
			
			\draw [ultra thick] (0,0) -- (\L,0) node [midway, below] {$m_1, L$};
			\draw [ultra thick] (0,0) -- (0,\L) node [midway, left] {$m_1, L$};
			
			\fill (\L,0) circle [radius = 2pt] node [above] {$m_2$};
			\fill (0,\L) circle [radius = 2pt] node [above] {$m_2$};
		\end{tikzpicture}
	\end{figure}
\end{example}

\begin{solution}
	\begin{align*}
		I &= \dfrac{2 m_1 L^2}{3} + 2 m_2 L^2
	\end{align*}
	~\\
	\begin{align*}
		x_{\textnormal{COM}} &= \dfrac{m_1 \cdot 0 + m_2 \cdot 0 + m_1 \cdot \sfrac{L}{2} + m_2 \cdot L}{2 m_1 + 2 m_2}\\
		\therefore x_{\textnormal{COM}} &= \dfrac{m_1 L + 2 m_2 L}{4(m_1 + m_2)}\\
		y_{\textnormal{COM}} &= \dfrac{m_1 L + 2 m_2 L}{4(m_1 + m_2)}\\
	\end{align*}
	~\\
	\begin{figure}[H]
		\begin{tikzpicture}
			\def\L{2};
			
			\draw [gray] ({-2*\L},0) -- ({2*\L},0);
			
			\draw [ultra thick] (0,0) -- (\L,0) node [midway, below] {$m_1, L$};
			\draw [ultra thick] (0,0) -- (0,\L) node [midway, left] {$m_1, L$};
			
			\fill (\L,0) circle [radius = 2pt] node [above] {$m_2$};
			\fill (0,\L) circle [radius = 2pt] node [above] {$m_2$};
		\end{tikzpicture}
		\caption{Stable equilibrium condition}
	\end{figure}	
	\begin{figure}[H]
		\begin{tikzpicture}
			\def\L{2};
	
			\draw [gray] ({-2*\L},0) -- ({2*\L},0);
			
			\draw [ultra thick] (0,0) -- (45:\L) node [midway, below right] {$m_1, L$};
			\draw [ultra thick] (0,0) -- (135:\L) node [midway, below left] {$m_1, L$};
			
			\fill (45:\L) circle [radius = 2pt] node [above] {$m_2$};
			\fill (135:\L) circle [radius = 2pt] node [above] {$m_2$};
		\end{tikzpicture}
		\caption{Unstable equilibrium condition}
	\end{figure}
	\begin{figure}[H]
		\begin{tikzpicture}
			\def\L{2};
			
			\draw [gray] ({-2*\L},0) -- ({2*\L},0);
			
			\draw [ultra thick] (0,0) -- (-\L,0) node [midway, below] {$m_1, L$};
			\draw [ultra thick] (0,0) -- (0,\L) node [midway, right] {$m_1, L$};
			
			\fill (-\L,0) circle [radius = 2pt] node [above] {$m_2$};
			\fill (0,\L) circle [radius = 2pt] node [above] {$m_2$};
		\end{tikzpicture}
		\caption{Stable equilibrium condition}
	\end{figure}
	For the body to topple,
	\begin{align*}
		\dfrac{1}{2} I {\omega_{\textnormal{min}}}^2 &= m g \left( \sqrt{{x_{\textnormal{COM}}}^2 + {y_{\textnormal{COM}}}^2} - y_{\textnormal{COM}} \right)
	\end{align*}
\end{solution}

\end{document}